\section{Arithmetic: estimating every coefficient at once}\label{sec:local-arithmetic}

The real integral makes the determinant positive and small. Arithmetic has
a different task: it must control \emph{all} coefficients of the polynomial,
including the constant coefficient contributed by polynomial parts in
partial fractions. Controlling only the residues is not sufficient.

This section follows Fauzan's choice of bases and local organization
\cite[\S\S3--4]{fauzan}, with the direct partial-fraction proofs, capped
inner weights, and full outer rank loss in Firsching's formalization
\cite[\texttt{Apery/Arith/} and \texttt{Norm.lean}]{firsching}.
These are source-derived constructions. We give their mathematical
proofs here rather than using the existence of formal code as a substitute
for the estimates.

Throughout this section,
\[
 K=40n,\qquad N=3n,\qquad h=37n,\qquad
 A_K(t)=D_N(t)^6/D_K(t),\qquad F_K=S_K\Delta_K.
\]
All valuation statements concern rational numbers, so no completion of
$\Q$ is needed. Write $\Z_{(p)}$ for rationals whose denominators are
prime to $p$, put $v_p(0)=+\infty$, and define
\[
 v_p^G\!\left(\sum_j a_jX^j\right)=\min_j v_p(a_j).
\]
Thus $v_p^G(P)\ge b$ means that \emph{each} coefficient has valuation
at least $b$. The superscript distinguishes this coefficient bound from
evaluation at a numerical value of $X$.

All three bases estimate the same form \(\mathcal B_{K,X}\) from
\eqref{eq:bilinear-form}. The small-prime basis exposes integer values,
the inner basis prescribes vanishing orders at residue classes, and
the outer basis cancels actual poles. Their coordinate costs are
accounted for by the following rule.

\subsection{Two elementary determinant rules}

If a matrix $B(X)$ satisfies
$v_p^G(B_{ij})\ge u_i+v_j$, then
\begin{equation}\label{arith:detweights}
 v_p^G(\det B)\ge\sum_i u_i+\sum_j v_j.
\end{equation}
Indeed, each permutation product has this lower bound, and addition
cannot lower it. The weights may be half-integers; the coefficients
themselves still have integer valuations.

If $E_0,\ldots,E_{h-1}$ is another basis of $\Q[t]_{<h}$ with
coefficient matrix $C$ defined by $E_i(t)=\sum_j C_{ji}t^j$,
bilinearity gives
\[
 \bigl[\mu_X(A_KE_iE_j)\bigr]=C^{\mathsf T}G_K(X)C.
\]
Consequently its determinant is $(\det C)^2\Delta_K(X)$.
When $C$ has entries in $\Z_{(p)}$ and determinant a unit there, the
valuation of the determinant is unchanged. Choosing such a basis
separately at each prime is the principal freedom in the argument.
These identities are the basis-change and weighted-determinant steps
used throughout \cite[\S\S3.3--4]{fauzan};
see also \cite[\texttt{BasisChange.lean}, \texttt{Val.lean}]{firsching}.

\subsection{The Bernoulli functional and integer-valued polynomials}

Let $\mathcal L(x^j)=B_j$ and
\[
 \tau(P)=\frac1{24}\mathcal L(P''').
\]
The same Bernoulli data supplied the real moments. With
\(\partial=d/dx\), their generating function gives the
Euler--Maclaurin operator
\[
 \tau(P)=
 \left[\frac{\partial^4}{24(e^\partial-1)}P\right](0).
\]
The series starts with \(\partial^3/24\) and terminates on each
polynomial \cite[24.2 and 24.4]{dlmf}. Since \(e^\partial-1\)
acts as forward difference, multiplication gives
\[
 \tau\bigl(P(x+1)-P(x)\bigr)=P^{(4)}(0)/24.
\]
This explains the use of finite differences below. The operator formula
is asserted only on polynomials.

For rational functions with simple integer poles, define
\[
 \tau_X\left(P(x)+\sum_r\frac{c_r}{x-r}\right)
 =\tau(P)+\sum_r c_r\bigl(H_{d(r)}^{(5)}-X\bigr),\qquad
 d(r)=
 \begin{cases}r,&r\ge0,\\-r-1,&r<0.\end{cases}
\]
The pullback identity, proved in Appendix~\ref{app:pullback}, is
\begin{equation}\label{arith:pullback}
 \mu_X(R(t))=\tau_X\bigl(x^5R(-x^2)\bigr).
\end{equation}
It turns square poles into integer poles. This identity and the
functional $\tau$ are Fauzan's, not new definitions introduced to repair
the proof \cite[\S3, (3.1)]{fauzan}.
For an odd prime, a class \(t=-c^2\) with \(c\ne0\) gives the two classes
\(x=\pm c\); at zero, vanishing orders double.

The useful basis for estimating $\tau$ is
$\binom{x}{j}=x(x-1)\cdots(x-j+1)/j!$.
For a polynomial $P$ of degree at most $d$, Newton's identity reads
\[
 P(x)=\sum_{j=0}^d(\Delta^jP)(0)\binom{x}{j},
 \qquad \Delta P(x)=P(x+1)-P(x).
\]
It follows by evaluating successively at $0,1,\ldots,d$, since the
evaluation matrix is triangular with diagonal one.
If $v_p(P(m))\ge\beta$ for $0\le m\le d$, then every coefficient
$(\Delta^jP)(0)$ has valuation at least $\beta$: it is an integer
linear combination of those values.

Two generating-function identities quantify the cost of differentiation
and of applying $\mathcal L$:
\begin{align}
 \frac{d}{dx}\binom{x}{j}
 &=\sum_{\ell=1}^j\frac{(-1)^{\ell-1}}{\ell}
                   \binom{x}{j-\ell},\label{arith:bin-deriv}\\
 \mathcal L\!\left(\binom{x}{j}\right)&=\frac{(-1)^j}{j+1}.
 \label{arith:bin-bernoulli}
\end{align}
For the first, differentiate
$(1+z)^x=\sum_j\binom{x}{j}z^j$ in $x$ and use the series for
$\log(1+z)$. For the second, apply $\mathcal L$ formally to
$(1+z)^x=e^{x\log(1+z)}$ and use
$\mathcal L(e^{ux})=u/(e^u-1)$.
Each fixed coefficient involves finitely many terms, so these are
identities of formal series, with no convergence issue.

Set $\Lambda=\lfloor\log_p(d+1)\rfloor$.
Equation~\eqref{arith:bin-deriv} loses at most $\Lambda$ in valuation
at each differentiation, and \eqref{arith:bin-bernoulli} loses at most
another $\Lambda$. Division by $24$ gives
\begin{equation}\label{arith:taupoly}
 v_p(\tau(P))\ge
 \beta-4\lfloor\log_p(d+1)\rfloor-v_p(24).
\end{equation}
The same statement holds if the hypotheses are known on
\emph{any} $d+1$ consecutive integers. Newton expansion based at the
left endpoint first shows the bound at every integer: the numbers
$\binom{m}{j}$ are integers for every integer $m$, including negative
$m$. Then apply the displayed estimate in its original form.
This proves the polynomial estimate (3.9) of
\cite[\S3.3]{fauzan}; the direct finite-difference implementation is
\cite[\texttt{BinomBasis.lean}, \texttt{TauBound.lean},
\texttt{SmallPrime.lean}]{firsching}.

\subsection{A bound valid at every prime}

We first prove a deliberately coarse bound. Its role is to handle
the small primes and any finite exceptional cases.

\begin{lemma}\label{arith:crude-tau}
Let $A(x)$ be integer-valued on $\Z$, with $\deg A\le d$ and
$d\ge2K\ge2$. Put
\[
 g(x)=\frac{(K!)^2A(x)}
             {\prod_{r\in\{-K,\ldots,-1,1,\ldots,K\}}(x-r)}.
\]
Then, for every prime $p$,
\[
 v_p^G(\tau_X(g))
 \ge-6\lfloor\log_p(d+1)\rfloor-v_p(24).
\]
\end{lemma}

\begin{proof}
Write $\Lambda=\lfloor\log_p(d+1)\rfloor$ and decompose
$g=P+\sum_r c_r/(x-r)$. For $r=\pm j$, direct multiplication of the
differences between the poles gives
\[
 \frac{(K!)^2}
 {\prod_{s\ne r}(r-s)}
 =\pm j\,\frac{(K!)^2}{(K-j)!(K+j)!}
 =\pm j\,\frac{\binom{2K}{K+j}}{\binom{2K}{K}}.
\]
For $0\le v\le u$, Legendre's formula implies
\[
 v_p\binom uv
 =\sum_{a\ge1}\left(
 \left\lfloor\frac u{p^a}\right\rfloor
 -\left\lfloor\frac v{p^a}\right\rfloor
 -\left\lfloor\frac{u-v}{p^a}\right\rfloor\right)
 \le\lfloor\log_p u\rfloor;
\]
each summand is zero or one and vanishes for $p^a>u$.
Thus $v_p(c_r)\ge-\Lambda$.

For an integer $m>K$,
\[
 \frac{(K!)^2}{\prod_r(m-r)}
 =\frac1{\binom{m-1}{K}\binom{m+K}{K}}.
\]
Take $m=K+1,\ldots,d-K+1$. Here $m+K\le d+1$, so
$v_p(g(m))\ge-2\Lambda$. Also $1\le m-r\le d+1$ for every pole,
giving $v_p(c_r/(m-r))\ge-2\Lambda$. Therefore
$v_p(P(m))\ge-2\Lambda$ on these $d-2K+1$ consecutive integers.
Since $\deg P\le d-2K$, the window form of
\eqref{arith:taupoly} yields
$v_p(\tau(P))\ge-6\Lambda-v_p(24)$.

Finally $d(r)\le K\le d+1$, so
$v_p^G(H_{d(r)}^{(5)}-X)\ge-5\Lambda$.
Each residue term in $\tau_X(g)$ has valuation at least $-6\Lambda$.
Adding the polynomial and pole contributions proves the result.
\end{proof}

This proof is the integer-window argument in
\cite[\texttt{SmallPrimeTau.lean}, \texttt{crude\_tauX}]{firsching},
replacing the local analytic discussion in Fauzan's Lemma~3.3.

Now use
\[
 q_0(t)=1,\qquad
 q_i(t)=\frac{(-1)^i\,2tD_{i-1}(t)}{(2i)!}\quad(i\ge1).
\]
Two product identities explain this otherwise rather special basis:
\begin{align*}
 q_i(-z^2)&=\binom{z+i}{2i}+\binom{z+i-1}{2i}\quad(i\ge1),\\
 \frac{D_N(-z^2)}{(N!)^2}
 &=\binom{N+z}{N}\binom{N-z}{N}.
\end{align*}
The identities follow by expanding the products on both sides. In
particular both expressions are integers for every integer $z$.
Put $c=(K!)^2/(N!)^{12}$. By the pullback formula,
$c\,\mu_X(A_Kq_iq_j)$ is the $\tau_X$-value in
Lemma~\ref{arith:crude-tau} with integer-valued numerator
\[
 (-1)^Kx^5
 \left(\frac{D_N(-x^2)}{(N!)^2}\right)^6
 q_i(-x^2)q_j(-x^2).
\]
Its degree is at most
$5+12N+4(h-1)=184n+1$, whose successor is at most $5K=200n$.
Every matrix entry consequently has valuation at least
$-6\lfloor\log_p(5K)\rfloor-v_p(24)$.

The coefficient matrix of the $q_i$ is triangular, with leading
coefficients $1,(-1)^i2/(2i)!$. Hence
\[
 \det\bigl[c\,\mu_X(A_Kq_iq_j)\bigr]
 =c^h\frac{4^{h-1}}{\prod_{i=1}^{h-1}((2i)!)^2}\Delta_K
 =F_K.
\]
We have proved
\begin{equation}\label{arith:small}
 \boxed{\quad v_p^G(F_K)\ge
 -6h\lfloor\log_p(5K)\rfloor-hv_p(24).\quad}
\end{equation}
This is Fauzan's (3.12), with its direct proof in
\cite[\texttt{SmallPrimeF.lean}, \texttt{small\_prime\_bound}]{firsching}.

\subsection{A direct bound from residue-class counts}

For larger primes a coarse common denominator wastes too much
divisibility. We now count zeros and poles in each residue class.
The following finite argument is the key replacement for a
$p$-adic extension of the functional.

\begin{lemma}\label{arith:direct}
Let $p\ge5$, and let
\[
 g(x)=\kappa\frac{\prod_{\zeta\in Z}(x-\zeta)}
                         {\prod_{r\in P}(x-r)},\qquad
 \kappa\in\Z_{(p)},
\]
where $Z$ is a multiset of integers and $P$ is a finite set of distinct
integers. For a residue class $c$, let $e_c$ and $\ell_c$ count its
zeros and poles. Assume that distinct poles congruent modulo $p$
are not congruent modulo $p^2$, that $d(r)<p^2$ for all poles, and that
the polynomial part has degree at most $d$ with $d+1<p^2$.
Then
\[
 v_p^G(\tau_X(g))\ge\min_c(e_c-\ell_c)-4.
\]
\end{lemma}

\begin{proof}
Write $g=P_0+\sum_r c_r/(x-r)$ and
$\beta=\min_c(e_c-\ell_c)$.
The residue formula gives
\begin{equation}\label{arith:residuecount}
 v_p(c_r)\ge e_c-\ell_c+1,\qquad c=r\bmod p.
\end{equation}
Indeed, each numerator factor in class $c$ contributes at least one
power of $p$; the $\ell_c-1$ denominator factors in that class
contribute exactly one each, and all other denominator factors are units.

We need a bound on the polynomial part as well. Fix an integer $m$
in class $c$, and write the pole denominator as $\Pi_c\Pi_o$,
where $\Pi_c$ contains the poles in $c$ and $\Pi_o$ the other poles.
If $A$ denotes the numerator, the partial fractions give
\[
 A=\Pi_c H+\Pi_o I,
\]
where
\[
 H=P_0\Pi_o+\sum_{r\notin c}c_r\frac{\Pi_o}{x-r},
 \qquad
 I=\sum_{r\in c}c_r\frac{\Pi_c}{x-r}.
\]
Substitute $x=m+pz$. The coefficients of $A(m+pz)$ are divisible
by $p^{e_c}$ in $\Z_{(p)}$.
By \eqref{arith:residuecount}, the same is true of $I(m+pz)$:
its residue supplies $e_c-\ell_c+1$ powers and its remaining
$\ell_c-1$ factors supply the rest. Also
\[
 \Pi_c(m+pz)=p^{\ell_c}
       \prod_{r\in c}\left(z-\frac{r-m}{p}\right)=p^{\ell_c}T(z),
\]
with $T$ monic and integral. It follows that
$T(z)\,p^{\ell_c-e_c}H(m+pz)$ has integral coefficients.
Division by a monic integral polynomial preserves integral
coefficients, so $p^{\ell_c-e_c}H(m+pz)$ is integral. At $z=0$
this gives $v_p(H(m))\ge e_c-\ell_c$.

Since $\Pi_o(m)$ is a unit, the identity
\[
 P_0(m)=\frac{H(m)}{\Pi_o(m)}
              -\sum_{r\notin c}\frac{c_r}{m-r}
\]
now proves $v_p(P_0(m))\ge\beta$. This holds for every integer $m$.
Estimate~\eqref{arith:taupoly}, with $d+1<p^2$ and $p\nmid24$,
therefore gives $v_p(\tau(P_0))\ge\beta-4$.
The harmonic bound $d(r)<p^2$ gives
$v_p^G(H_{d(r)}^{(5)}-X)\ge-5$; together with
\eqref{arith:residuecount} it gives the same $\beta-4$ bound
for every pole term.
\end{proof}

This is the argument of
\cite[\texttt{DirectBound.lean}, \texttt{DirectPoly.lean},
\texttt{tauX\_direct}]{firsching}. The count measures the whole
rational function; the polynomial-part paragraph is essential.

\subsection{Inner primes: distributing vanishing among the classes}

Write $p=2m+1$ and represent classes up to sign by
$c=0,1,\ldots,m$. Put
\begin{align*}
 s_H(c)=\#\{1\le j\le H:j\equiv c\pmod p\}
       &+\#\{1\le j\le H:j\equiv-c\pmod p\},\\
 \beta_c&=6s_N(c)-s_K(c)-4.
\end{align*}
The two counts coincide when $c=0$, so $s_H(0)=2\lfloor H/p\rfloor$.
This doubled convention is intentional.

Choose nonnegative integers $L_c$ with $\sum_cL_c=h$ and define
\[
 E_{c,i}(t)=(t+c^2)^i\prod_{a\ne c}(t+a^2)^{L_a},
 \qquad 0\le i<L_c.
\]
These form a basis with coefficient determinant a $p$-adic unit.
To see this, reduce modulo $p$ and work in
$\mathbf F_p[t]/\prod_c(t+c^2)^{L_c}$.
The factors for different $c$ are coprime. In the $c$-component,
the rows belonging to $c$ are $(t+c^2)^i$ times a unit, while
all other rows vanish. They are therefore a basis in each component,
and the Chinese remainder theorem proves the assertion.
This auxiliary modulus prescribes Taylor coefficients at each class;
it need not equal the pole denominator \(D_K/D_N\). The allocation
\(L_c\) distributes the same \(h\) polynomial coordinates among classes.

For a row $u=(c,i)$ let
$\nu_u(a)=i$ if $a=c$ and $\nu_u(a)=L_a$ otherwise.
In the pullback of $A_KE_uE_v$, the count of zeros minus poles in
either residue class represented by $a$ is
\[
 \begin{cases}
 6s_N(a)-s_K(a)+\nu_u(a)+\nu_v(a),&a\ne0,\\
 5+6s_N(0)-s_K(0)+2\nu_u(0)+2\nu_v(0),&a=0.
 \end{cases}
\]
The extra $5$ comes from $x^5$; the factor $2$ comes from
$t=-x^2$ at the zero class.
All hypotheses of Lemma~\ref{arith:direct} hold, for example, if
$p^2>3K$: distinct integer poles lie between $-K$ and $K$,
the harmonic indices are at most $K$, and the polynomial-part
degree plus one is at most $104n+2<3K$.

Set
\[
 T=\min_{a\ne0}\left(L_a+\frac{\beta_a}{2}\right)
\]
and assign row weights
\begin{equation}\label{arith:innerweights}
 w_{c,i}=
 \begin{cases}
 \min\{i+\beta_c/2,T\},&c\ne0,\\
 \min\{2i+(\beta_0+5)/2,T\},&c=0.
 \end{cases}
\end{equation}
Assume the zero-class condition
\begin{equation}\label{arith:zero-condition}
 T\le2L_0+\frac{\beta_0+5}{2}.
\end{equation}
For a nonzero class $a$, a row belonging to $a$ has weight at most
$i+\beta_a/2$; any other row has weight at most
$T\le L_a+\beta_a/2$. At $a=0$ the analogous bounds follow from
\eqref{arith:zero-condition}. Adding the bounds for the two rows
shows, class by class, that the minimum in
Lemma~\ref{arith:direct}, less $4$, is at least $w_u+w_v$.
The basis-change and determinant rules therefore give
\begin{equation}\label{arith:inner}
 \boxed{\quad v_p^G(\Delta_K)\ge2\sum_{c=0}^m\sum_{i=0}^{L_c-1}w_{c,i}.\quad}
\end{equation}
This is the capped-weight bound proved as \texttt{inner\_bound\_cap}
in \cite[\texttt{InnerWeights.lean}]{firsching}. Its basis is
Fauzan's class basis \cite[\S4.1]{fauzan}.

\subsection{Choosing the inner allocation}

The allocation can be specified without an optimization oracle.
Put $q=\lfloor K/p\rfloor$ and
\[
 k_{\rm lo}=-2q-6,\qquad k_{\rm top}=8q+20,\qquad L_{\rm reserve}=3q+6.
\]
For an integer $k$ define
\[
 \ell(k,\beta)=\max\left(0,\left\lceil\frac{k-\beta}{2}\right\rceil\right),
 \qquad R(k)=\sum_{c=1}^m\ell(k,\beta_c).
\]
The number $\ell(k,\beta_c)$ counts the row weights
$i+\beta_c/2$ that are strictly below $k/2$.
Choose the largest integer $k_*$ in $[k_{\rm lo},k_{\rm top}]$
such that $R(k_*)\le h-L_{\rm reserve}$, and set
\[
 L_c=\ell(k_*,\beta_c)\ (c\ne0),\qquad L_0=h-R(k_*).
\]
For sufficiently large $K$, uniformly when $K/400<p\le K/3$,
this choice exists. In fact $\beta_c\ge-2q-6=k_{\rm lo}$, so
$R(k_{\rm lo})=0$, and $L_{\rm reserve}\le3\cdot399+6<h$.
All parameters depending on $q$ are bounded absolutely in this range.

The zero-class condition also follows directly. Since $N\le K$,
\[
 -2q-4\le\beta_0,\qquad
 \beta_c\le10q+8\quad(c\ne0).
\]
The ceiling formula implies
$L_c+\beta_c/2\le\max(\beta_c,k_*+1)/2$, whence
$T\le5q+21/2$.
On the other hand,
$2L_0+(\beta_0+5)/2\ge5q+25/2$.
Thus \eqref{arith:zero-condition} holds with room to spare.

Here is the useful quantitative interpretation. Define
\[
 \psi(k,\beta)=\sum_{i\ge0}\frac{(k-\beta-2i)_+}{2},
 \qquad
 \Phi(k)=\frac{hk}{2}-\sum_{c=1}^m\psi(k,\beta_c),
 \qquad x_+=\max(x,0).
\]
Each sum defining $\psi$ is finite. Ordinary rows are not truncated
by the cap: their raw weights are below $k_*/2$, whereas $T\ge k_*/2$.
The zero rows contribute at least $L_0k_*/2-B$, where
\[
 B\le \sum_{i\ge0}
 \left(\frac{k_*-\beta_0-5}{2}-2i\right)_+.
\]
Because $q$, $k_*$ and $\beta_0$ are bounded in the inner range,
$B=O_{400}(1)$, independently of $K$ and $p$. Consequently
\[
 \sum_{c,i}w_{c,i}\ge\Phi(k_*)-O_{400}(1).
\]
Finally,
\[
 \Phi(k+1)-\Phi(k)=\frac{h-R(k+1)}2.
\]
For $k<k_*$ the increments are nonnegative. After $k_*$,
maximality implies $R(k+1)>h-L_{\rm reserve}$, so each increment
is less than $L_{\rm reserve}/2$. Therefore every integer
$k\in[k_{\rm lo},k_{\rm top}]$ satisfies
\begin{equation}\label{arith:inner-comparison}
 2\sum_{c,i}w_{c,i}
 \ge hk-2\sum_{c=1}^m\psi(k,\beta_c)-O_{400}(1).
\end{equation}
The constant is uniform: one may bound the additional loss by
$(k-k_{\rm lo})L_{\rm reserve}+2B$, and every factor here is
bounded in terms of $400$.

This is the water-filling argument of
\cite[\texttt{InnerAlloc.lean}, \texttt{PhiK\_le\_kf}]{firsching}.
It explains why the global calculation may try several convenient
levels $k$, although the normalizing factor itself uses one fixed,
explicit allocation.

\subsection{Outer primes: separate the large poles from a small correction}

Assume
\[
 p\ge7,\qquad p>K/3,\qquad p^2>2K.
\]
Then $N<p/2$ and $K<3p$. The class basis now uses the actual remaining
poles. For $c=0,\ldots,m$ put
\[
 J_c=\{j:N<j\le K,\ j\equiv\pm c\pmod p\},\quad
 L_c=|J_c|,\quad
 B_c=\{j\in J_c:j>p\},\quad b_c=|B_c|.
\]
Define
\begin{align*}
 R_c(t)&=\prod_{\substack{N<j\le K\\j\notin J_c}}(t+j^2),\\
 q_{c,i}(t)&=
 \begin{cases}
 (t+c^2)^i,&i<b_c,\\
 \displaystyle\prod_{j\in B_c}(t+j^2)(t+c^2)^{i-b_c},&i\ge b_c,
 \end{cases}\\
 E_{c,i}(t)&=R_c(t)q_{c,i}(t),\qquad 0\le i<L_c.
\end{align*}
Modulo $p$, these have the same Chinese-remainder form as the inner
basis, so their coefficient determinant is again a unit.
The distinction between $i<b_c$ and $i\ge b_c$ has an exact purpose:
the latter rows cancel every pole $j>p$ in their own class.
This basis is Fauzan's \cite[\S4.2]{fauzan}; its precise use below follows
\cite[\texttt{OuterBasis.lean}, \texttt{OuterEntryMain.lean}]{firsching}.

Let
\[
 \ell_c=\#\{1\le j\le K:j\equiv\pm c\pmod p\},
 \qquad \delta_c=\mathbf1_{\{c\le N\}}\quad(c\ne0).
\]
Here, unlike $s_H(0)$ above, $\ell_0$ counts each multiple of $p$
only once. Use the following nonpositive weights:
\begin{equation}\label{arith:outerweights}
 w_{c,i}=
 \begin{cases}
 \min\{0,i+3\delta_c-(\ell_c+4)/2\},&c\ne0,\ i<b_c,\\
 0,&c\ne0,\ i\ge b_c,\\
 -2,&c=0,\ i<b_0,\\
 -1/2,&c=0,\ b_0=0,\\
 0,&c=0,\ i\ge b_0>0.
 \end{cases}
\end{equation}
The cases concern existing rows $i<L_c$; empty classes contribute nothing.
Since $K<3p$, the zero class has $b_0=0$ or $1$.

\subsubsection*{The polynomial moments that can fail to be integral}

The von Staudt--Clausen theorem states, for $r\ge1$, that
$B_{2r}+\sum_{q-1\mid2r}1/q$ is an integer, the sum being over
primes $q$ \cite[24.10(i)]{dlmf}.
Thus $v_p(B_{2r})\ge-1$ and it is nonnegative unless
$p-1\mid2r$. For our moments
\[
 \mu(t^e)=\frac{(-1)^eB_{2e+2}(2e+3)(2e+4)(2e+5)}{24},
\]
this gives $v_p(\mu(t^e))\ge-1$. In fact
\begin{equation}\label{arith:low-moments}
 v_p(\mu(t^e))\ge0\qquad(e<2p-3).
\end{equation}
To verify the stronger assertion, if $p-1\mid2e+2$ in this range,
write $2e+2=j(p-1)$. Necessarily $j=1,2,$ or $3$, and
$2e+2+j=jp$ is one of the three consecutive factors in the numerator.
It cancels the possible denominator $p$.

For an integral polynomial $Q$, let $\mu_{\ge}(Q)$ be the sum of
the terms of $\mu(Q)$ arising from its monomials of degree at least
$2p-3$. Then
$\mu(Q)-\mu_{\ge}(Q)\in\Z_{(p)}$ and
$p\mu_{\ge}(Q)\in\Z_{(p)}$.
For each matrix entry, subtract $\mu_{\ge}$ of its polynomial part,
and denote the resulting matrix by $A(X)$. We shall prove
\begin{equation}\label{arith:outer-entry}
 v_p^G(A_{uv})\ge w_u+w_v.
\end{equation}
The omitted polynomial part will then be restored with its full
rank cost. The moment argument is
\cite[\texttt{MomentVal.lean}, \texttt{OuterFrame.lean}]{firsching}.

\subsubsection*{Pole estimates, including the cancellation between two poles}

Write
$\varphi_j(X)=j^4(X-H_j^{(5)})-\tfrac14+\tfrac1{2j}$.
Directly from the denominators in the harmonic sum,
\[
 v_p^G(\varphi_j)\ge
 \begin{cases}
 0,&1\le j<p,\\
 -5,&1\le j<p^2,\\
 -1,&j=pu,\ 1\le u<p.
 \end{cases}
\]
In the last case the factor $j^4$ recovers four of the five possible
powers lost in the harmonic denominator.
We also need
\begin{equation}\label{arith:pole-congruence}
 \varphi_{p-a}(X)\equiv\varphi_a(X)\pmod p
 \quad(1\le a<p).
\end{equation}
Indeed, pairing $u$ with $p-u$ gives $H_{p-1}^{(5)}\equiv0$.
Thus $H_{p-a}^{(5)}\equiv H_{a-1}^{(5)}\pmod p$.
Now $(p-a)^4\equiv a^4$, and the difference in the harmonic terms
contributes $a^{-1}$, exactly canceled by
$1/(2(p-a))-1/(2a)\equiv-a^{-1}$.
These are the pole estimates of
\cite[\texttt{PoleVal.lean}]{firsching}.

Consider the rational function $A_KE_{c,i}E_{d,j}$.
If $c\ne d$, every denominator factor is canceled by the numerator,
so its pole sum is zero. Its polynomial part has integer coefficients
because division by the monic $D_K$ preserves integrality.
After removal of the high moments, \eqref{arith:low-moments}
therefore proves \eqref{arith:outer-entry} for different classes.

For a common nonzero class $c$, let $r_s$ be the residue at $t=-s^2$.
Since $p^2>2K$, for $p\nmid s$ each denominator factor
$u^2-s^2=(u-s)(u+s)$ has valuation one when $u\equiv\pm s$,
and zero otherwise. It cannot have two powers of $p$:
$|u-s|,u+s<p^2$, and $p$ cannot divide both $u-s$ and $u+s$.
The numerator has at least $6\delta_c+i+j$ such factors, hence
\begin{equation}\label{arith:outer-residue}
 v_p(r_s)\ge6\delta_c+i+j-\ell_c+1.
\end{equation}
If both $i,j<b_c$, adding the bound $-5$ for $\varphi_s$ gives
$6\delta_c+i+j-\ell_c-4$, which is at least the sum of the
two weights in \eqref{arith:outerweights}.

Suppose instead that $i\ge b_c$. All poles above $p$ disappear.
At most $c$ and $p-c$ remain. If $\delta_c=1$, the six numerator
factors in \eqref{arith:outer-residue} make each remaining residue
integral: $\ell_c\le b_c+2$ and $i\ge b_c$, so the lower bound
is at least $5$. If $\delta_c=0$ and at most one small pole occurs, then
$\ell_c\le b_c+1$, so the same is true. The only delicate case
has $\delta_c=0$ and both small poles $a=c$, $b=p-c$.
Then \eqref{arith:outer-residue} gives $v_p(r_a),v_p(r_b)\ge-1$.

All other denominator factors cancel exactly. Thus the rational
function is $S(t)/((t+a^2)(t+b^2))$ for an integral polynomial $S$,
and
\[
 r_a=\frac{S(-a^2)}{b^2-a^2},\qquad
 r_b=\frac{S(-b^2)}{a^2-b^2}.
\]
The sum $r_a+r_b$ is an integral divided difference of $S$.
Using \eqref{arith:pole-congruence},
\[
 r_a\varphi_a+r_b\varphi_b
 =(r_a+r_b)\varphi_a+r_b(\varphi_b-\varphi_a)
\]
is integral too. This is the needed cancellation; it is stronger
than separate bounds on the two residues.
It proves \eqref{arith:outer-entry} whenever one row has $i\ge b_c$,
because that row has weight zero and every other weight is nonpositive.
The argument is the pairing lemma and mixed case in
\cite[\texttt{OuterEntries.lean}, \texttt{pair\_GV};
\texttt{OuterEntryMain.lean}, \texttt{pole\_case\_b}]{firsching}.

In the zero class the poles are among $p,2p$. A factor $u^2-s^2$
between two distinct multiples of $p$ has valuation exactly two.
The numerator has at least $i+j$ such factors, and there are at most
$b_0$ denominator factors of this kind besides the pole under
consideration. Therefore
\[
 v_p(r_s\varphi_s)\ge2(i+j)-2b_0-1.
\]
For $b_0=0$ this is at least the sum of the two weights $-1/2$.
For $b_0=1$, the weights are $-2$ at $i=0$ and $0$ thereafter;
checking whether zero, one or both indices vanish proves the same
inequality. This completes \eqref{arith:outer-entry}.

\subsubsection*{Restoring the high moments and paying the rank cost}

In the original monomial basis, let
\[
 B_{ij}=\mu_{\ge}\left(
       \text{polynomial part of }D_N(t)^6t^{i+j}/D_K(t)\right).
\]
Its polynomial part has degree at most $i+j+6N-K=i+j-22n$.
Consequently $B_{ij}=0$ whenever $i+j<2p+22n-3$.
In particular the first
$k_0=\max(0,2p-15n-2)$ rows vanish, or all rows if $k_0\ge h$.
There are therefore at most
\[
 r_p=\max(0,h-k_0)=\max(0,52n+2-2p)
\]
remaining rows in the outer range.
Also $pB$ is integral. After the integral change of basis the
correction consequently has a factorization $p^{-1}UV$, where
$U$ is $h$ by $r_p$, $V$ is $r_p$ by $h$, and both are integral.
For example, factor $B$ through the projection onto its last
$r_p$ rows, and absorb the two basis-change matrices into the factors.

For completeness the determinant bound used here is
\[
 v_p^G\det(A+p^{-1}UV)\ge2\sum_u w_u-r_p
 \quad\text{if }w_u\le0,\quad v_p^G(A_{uv})\ge w_u+w_v.
\]
Indeed,
\[
 \det(A+p^{-1}UV)=
 \det\begin{pmatrix}A&-U\\p^{-1}V&I_{r_p}\end{pmatrix}.
\]
Apply \eqref{arith:detweights} with row weights $(w_u,-1,\ldots,-1)$
and column weights $(w_u,0,\ldots,0)$. Integrality of $U,V$ and
$w_u\le0$ check each block immediately.
This proves
\begin{equation}\label{arith:outer}
 \boxed{\quad
 v_p^G(\Delta_K)\ge
 2\sum_{c=0}^m\sum_{i=0}^{L_c-1}w_{c,i}-r_p.\quad}
\end{equation}
The full loss $r_p$ is the one used in
\cite[\texttt{LowRank.lean}, \texttt{OuterFrame.lean},
\texttt{OuterEntryMain.lean}]{firsching}. We do not replace it by
the stronger, disputed loss in the original manuscript.

\subsection{The local bounds are uniform}

Equations \eqref{arith:small}, \eqref{arith:inner}, and
\eqref{arith:outer} concern the same determinant and its normalization.
The inner and outer arguments apply uniformly for $p>K/400$ once
$K>480000$: then $p^2>3K$, the primes are at least $7$, and
$h>1203$ accommodates all reserved zero-class rows.
For smaller $K$ the coarse bound remains valid at every prime.
Thus a finite number of exceptional indices can be discarded, or handled
by the coarse bound, without changing any asymptotic argument.
The next section uses these explicitly defined weights to construct
the multiplier and bound its growth.
